\section{K-028: prime-to-three stabilizers and a smaller exceptional support}

Throughout this section $n\ge4$, $M=2^{n-1}$, and $\tau=(123)$.
Use $X_n,Y=X_n^{\mu_M},Y_n,S,\overline G$, and $\chi$ as in
K-026--K-027.  For a nonzero quotient point $y$, write $N_y$ for
the stabilizer of its closed orbit in the effective special-linear
group $\ker(\chi:\overline G\to\mathbb G_m)$.
Let $\Omega\subset X_n$ be the set where $N_y$ is finite and
its order is prime to three.

\begin{proposition}[K-028a]
The set $\Omega$ is open.  Let $W\subset Y_{\rm red}$ be an
irreducible locally closed subvariety with $\tau W=W$ and with
a dense open subset in $\Omega$.  Then $W\subseteq Y_n$.
In particular, a nonzero $\tau$-fixed point outside the chain
has either a positive-dimensional effective special-linear
stabilizer or a finite stabilizer whose order is divisible by
three.
\end{proposition}

\begin{proof}
At a closed orbit with finite stabilizer, Luna's slice theorem
gives an open saturated neighborhood in which every stabilizer
is conjugate to a subgroup of that finite stabilizer.  Its
quotient neighborhood therefore lies in $\Omega$ when the original
order is prime to three.  This proves openness.  The precise
primary statement is Luna (1973), Section III.1, slice theorem
and proof on printed p.~97, followed by Remark~$4^\circ$ on
printed p.~98; the latter explicitly gives the subgroup assertion.

Fix a primitive $M$th root $\xi$.  Shrink $W$ to a dense open
subset, invariant under $\tau$, where $|N_y|=q$ is constant and
$3\nmid q$.  Stabilizer orders and the finite counts below are
constructible, so such a shrinking is legitimate.  Here is a
direct way to see the particular constructibility needed.  Over
the stable locus in $V_n$, the stabilizer incidence scheme has
finite reduced fibres.  The cardinalities of fibres of a
quasi-finite morphism are constructible.  The same is true after
intersection with any fixed semisimple conjugacy class in
$\overline G$.  These counts are unchanged on an effective
special-linear orbit, and Chevalley's theorem therefore gives
constructible level sets on the quotient.  There are only
finitely many values relevant after $q$ is fixed.

Choose a closed tensor representative of $y$ and let $H_y$ be
its full product stabilizer modulo $S$.  It is finite: its
kernel under $\chi$ is $N_y$, and a nonzero homogeneous invariant
forces the image of $\chi$ to be finite.  Since $y\in Y$, the
coset
\[
                 E_y=\{h\in H_y:\chi(h)=\xi\}
                                                        \tag{K-028a}
\]
is nonempty and has exactly $q$ elements.  Every $h\in E_y$
satisfies $h^M\in N_y$, hence $h^{Mq}=1$.  Put $L=Mq$;
its order is prime to three.

There are finitely many $\overline G$ conjugacy classes of
elements of order dividing $L$.  For each such class $\mathcal C$,
let $m_{\mathcal C}(y)$ be the number of elements of $E_y$ in
that class.  This count is independent of the tensor representative.
After another dense open shrinking all these counts are constant.
Because $\tau W=W$, cyclic factor permutation gives
\[
 m_{\mathcal C}=m_{\tau\mathcal C},\qquad
                     \sum_{\mathcal C}m_{\mathcal C}=q.
                                                        \tag{K-028b}
\]
If every class with nonzero multiplicity had an orbit of size
three, the sum would be divisible by three.  Thus a
$\tau$-stable class occurs with nonzero multiplicity.

Choose an element of that class fixing a representative of a
point of the dense open subset.  As in K-026, lift it to $GL_n^3$
and take $L$th roots in the scalar torus $S$ to make the lift
have $L$th power one.  Its three spectral multisets are scalar
translates of one another.  Since $3$ is invertible modulo $L$,
scalar shifts of total zero make them equal.  After independent
conjugation in the three factors the fixing element is
$(A,A,A)$ with $A^L=I$ and $(\det A)^3=\xi$.

Take the two-primary part $A_2$ of this finite-order matrix.
The triple $(A_2,A_2,A_2)$ is a power of $(A,A,A)$ and hence
still fixes the tensor.  Its determinant has order exactly $M$:
cubing is injective on the finite cyclic group of order dividing
$L$, and the two-primary part of $\xi$ has order $M$.
This step discards any odd factors of $q$ other than three;
none are silently excluded from the original stabilizer.

The support-folding argument of K-026 now applies without any
freeness hypothesis.  Product-torus semistability gives cubic
support with $\mathbf1$ in its rational span.  The full support
theorem P-F30 at determinant order $M$ forces the complete cubic
chain support, and the unique positive tensor margins force
all ordered chain coefficients to be nonzero.  Diagonal-torus
transitivity puts every such tensor in the chain's product-linear
orbit.  A dense subset of $W$ is therefore contained in $Y_n$.
As $Y_n$ is closed, all of $W$ is contained in it.  Applying this
to the one-point subvariety proves the final assertion.
\end{proof}

\begin{corollary}[K-028b: refined exceptional class]
In K-027 one can replace $Z$ by the smaller closed set
\[
 Z_{(3)}=(Y_{\rm red}\setminus\Omega)
       \ \cup\!
       \bigcup_{\substack{Y_i\ne Y_n\ \text{an irreducible component}\\
                    Y_i\cap X_n^{\mu_{3M}}\not\subseteq\{0\}}}Y_i.
                                                        \tag{K-028c}
\]
Thus an integral virtual coherent class $\beta_{(3)}$ supported
on $Z_{(3)}$ satisfies exactly the Laurent congruence (K-027c)
and its necessary principal-part conditions for cancellation.
Finite stabilizers of order prime to three do not need to be
included in the exceptional support solely because they are
nontrivial.
\end{corollary}

\begin{proof}
The set $Z_{(3)}$ is closed, grading-stable, and $S_3$-stable.
It contains every component whose generic point is outside
$\Omega$, as well as every component closure with a nonzero
$\mu_{3M}$-fixed point.  It contains the origin but no nonzero
point of the chain.  The latter assertion uses P-F60: the chain
has trivial effective special-linear stabilizer, exact grading
stabilizer $\mu_M$, and a smooth one-dimensional fixed germ, so
it cannot meet another fixed component away from the origin.

Use the same integral Koszul class $\alpha$ as in K-027, retaining
its $\langle\tau\rangle$-equivariant structure.  The coherent
localization argument also holds in the category equivariant for
$\mathbb G_m\times\langle\tau\rangle$: when extending a coherent
sheaf, take the sum of the finitely many translates of the chosen
graded coherent submodule of its equivariant pushforward.  Morphisms
and exact-sequence relations can be extended equivariantly in the
same manner.  Thus the remainder classes at every step below may
be taken in this equivariant coherent group.  No division of a
Grothendieck class by three is used.  On
$Y\setminus Z_{(3)}$, the chain orbit is an open-and-closed
piece.  Extend its restriction equivariantly to a class $\gamma$ on the chain
closure as before.  It remains to treat the complement of that
piece.  This complement need not be smooth, and its components
need not be disjoint.  We use finite dimension induction instead
of a smoothness assertion.

Every irreducible component of that complement is permuted in
a three-element orbit by $\tau$, by K-028a.  Remove their pairwise
intersections.  The resulting dense open pieces are disjoint
and permuted in triples.  Choose coherent extensions from one
piece in each triple, supported on its closure, and take the
induced equivariant class on its three translates.  Its restriction
agrees, equivariantly, with the corresponding
restrictions of $\alpha-\gamma$.  Coherent localization leaves
a class supported on a proper closed subset.  Its intersection
with $Y\setminus Z_{(3)}$ has smaller dimension.  Every
$\tau$-stable irreducible component of this residual subset
would also contradict K-028a, since it lies outside the chain
and inside $\Omega$.  Repeating the operation therefore yields,
after finitely many dimension drops, an integral decomposition
\[
 \alpha=\gamma+
       \sum_i(\eta_i+\tau\eta_i+\tau^2\eta_i)
       +\beta_{(3)},
          \qquad \operatorname{Supp}\beta_{(3)}\subseteq Z_{(3)}.
                                                        \tag{K-028d}
\]
All supports here can be chosen grading-stable, by using homogeneous
ideals, grading-equivariant restrictions, and coherent extensions.

Every closure supporting an $\eta_i$ is contained in an original
fixed component not removed in (K-028c).  Hence it has no nonzero
$\mu_{3M}$-fixed point.  This remains true for all lower-dimensional
intersection closures introduced by the induction.  The complete
no-collision integrality lemma of K-027 applies to each such
Hilbert term.  Their translates have identical Hilbert series,
so the sum contributes a multiple of three in
$\mathcal O((1-t^M))$.  The chain contribution and the scalar
denominator $D_M$ are unchanged.  The remainder of the proof
of K-027 now gives its congruence and all its principal-part
consequences with $\beta_{(3)}$ in place of $\beta$.
\end{proof}

\paragraph{Primary locator and scope.}
Luna, D. (1973). \emph{Slices \'etales}. In \emph{Sur les groupes
alg\'ebriques}, \emph{Bulletin de la Soci\'et\'e Math\'ematique de
France, M\'emoire 33}, 81--105. doi:10.24033/msmf.110.
The complete Section III.1 slice theorem and proof, its preceding
lemma, and Remarks $1^\circ$--$4^\circ$ on printed pp.~96--98 were
read for this application; the finite locally closed stratification
is also stated in Section III.2, Corollaries 3--4, printed p.~100.
The reduction does not show that $Z_{(3)}$ has no positive-dimensional
parts and does not determine the exceptional class's Hilbert residue.
In particular, it is not an all-$n$ exact-$M$ pole theorem.
